% answers_short.tex —— 2016 年·全国Ⅰ卷·理科 · 简短答案速查
% 由 tools/make_short_answers.py 自 answers.tex 的 \gksolhead/\gkanswer 提取，
% 并与 2016/2016_数学真题_全国1理_答案要点.md 逐题比对后生成，勿手改（改源 MD 后重跑）。

\begin{gknotice}
\textbf{使用说明}：本节只给最终答案便于快速核对，完整推导见解析版。答案为编写者独立推导，\textbf{非官方标准答案}。
\end{gknotice}

\gkpart{选择题}
\begin{center}\begin{tabular}{ccccccccccccc}
\toprule
题号 & 1 & 2 & 3 & 4 & 5 & 6 & 7 & 8 & 9 & 10 & 11 & 12 \\
\midrule
答案 & D & B & C & B & A & A & B & C & C & B & A & B \\
\bottomrule
\end{tabular}\end{center}

\gkpart{填空题}
\begin{itemize}[leftmargin=2.6em]
  \item[\textbf{13.}] \(-2\)
  \item[\textbf{14.}] \(10\)
  \item[\textbf{15.}] \(64\)（\(n=3,4\) 时均取到）
  \item[\textbf{16.}] \(216000\) 元
\end{itemize}

\gkpart{解答题}
\begin{itemize}[leftmargin=2.6em]
  \item[\textbf{17.}] (1) \(C=60^\circ\)；(2) 周长 \(5+\sqrt7\)
  \item[\textbf{18.}] (1) 平面 \(ABEF\perp\) 平面 \(EFDC\)；(2) \(\cos\angle(E-BC-A)=-\frac{2}{\sqrt{19}}\)
  \item[\textbf{19.}] (1) \(P(X=x)\) 对应 \(x=16,17,18,19,20,21,22\) 依次为 \(0.04,0.16,0.24,0.24,0.20,0.08,0.04\)；(2) \(n_{\min}=19\)；(3) 购买 \(19\) 个；期望费用为 \(4040\) 元（购买 \(20\) 个时为 \(4080\) 元）
  \item[\textbf{20.}] (1) \(|EA|+|EB|=4\)；轨迹为 \(\frac{x^2}{4}+\frac{y^2}{3}=1\)，但不含 \((\pm2,0)\)；(2) 面积范围 \([12,8\sqrt3)\)
  \item[\textbf{21.}] (1) \(a>0\)；(2) \(x_1+x_2<2\)
\end{itemize}

\gkpart{选做题}
\begin{itemize}[leftmargin=2.6em]
  \item[\textbf{22.}] (1) \(AB\) 与圆相切；(2) \(AB\parallel CD\)
  \item[\textbf{23.}] (1) 圆心 \((0,1)\)、半径 \(a\) 的圆；\(\rho^2-2\rho\sin\theta+1-a^2=0\)；(2) U；若另加 \(C_1\cap C_2\ne\varnothing\)，则 \(a=1\)。见 ISSUE-0933
  \item[\textbf{24.}] (1) \(f(x)=\begin{cases}x-4,\&x<-1,\\3x-2,\&-1\le x<\frac32,\\4-x,\&x\ge\frac32\end{cases}\)；(2) \((-\infty,\frac13)\cup(1,3)\cup(5,+\infty)\)
\end{itemize}
